\section{Choosing disjoint auxiliary primes on the common window}
\label{sec:level-prime-selection}

We give the elementary part of the prime selection used by
Section~\ref{sec:host-field}, retaining its modulus \(N\), prime \(d\),
rank \(r\), and scale \(x\).
For \(g\geq2\) and \(r\geq1\) define
\[
 {\cal B}_r(g)=\{g\}\cup
 \bigcup_{j=1}^{r}
 \{\,\ell\text{ prime}:g^j<\ell^{4v_\ell(g^j-1)}\,\}.
\]
These sets are the \srcid{BB} of the source. We first prove the
count which makes two such sets disjoint.

\subsection{Counting exceptional divisors}

For fixed \(g,j\), at most three primes occur in the \(j\)-th set.
Every such prime divides the positive integer \(n=g^j-1\).
If four distinct primes \(\ell_i\) occurred, their powers
\(\ell_i^{v_{\ell_i}(n)}\) would be pairwise coprime divisors of \(n\),
so
\[
 (g^j)^4<
 \prod_{i=1}^{4}\ell_i^{4v_{\ell_i}(n)}
 =\left(\prod_{i=1}^{4}\ell_i^{v_{\ell_i}(n)}\right)^4
 \leq n^4<(g^j)^4 .
\]
This is impossible. Each set is finite because all its elements
divide \(n\). Consequently
\begin{equation}
 \label{eq:level-bb-card}
       |{\cal B}_r(g)|\leq3r+1 .
\end{equation}

We also need the bound
\begin{equation}
 \label{eq:level-roots}
 \#\{a\bmod\ell^w:a^j=1\}\leq2j
       \qquad(w,j\geq1).
\end{equation}
Such residues are units. For odd \(\ell\) the unit group modulo
\(\ell^w\) is cyclic. Here is the usual lifting justification.
The multiplicative group of the residue field is cyclic: if \(m\)
is its exponent, the polynomial \(X^m-1\) has at most \(m\) roots,
while all its elements are roots. In a finite abelian group the
exponent is the order of an element, by taking commuting elements
with maximal order in each primary component. Thus \(m\) equals
the group order. Choose a primitive residue \(a\bmod\ell\).
One of the \(\ell\) lifts \(a+k\ell\bmod\ell^2\) satisfies
\[
          v_\ell((a+k\ell)^{\ell-1}-1)=1 ,
\]
because modulo \(\ell^2\) the coefficient of \(k\) is
\((\ell-1)a^{\ell-2}\ell\), which is nonzero after division by
\(\ell\). If \(z=1+\ell^b u\), \(\ell\nmid u\), \(b\geq1\),
the binomial expansion for odd \(\ell\) gives
\(v_\ell(z^\ell-1)=b+1\). Induction shows that this lift has order
\((\ell-1)\ell^{w-1}\) modulo \(\ell^w\), proving cyclicity.
In a cyclic group of order \(m\), the equation \(a^j=1\) has
\(\gcd(j,m)\leq j\) solutions.

For \(\ell=2\), the moduli \(2\) and \(4\) satisfy
\eqref{eq:level-roots} directly. When \(w\geq3\), repeated squaring
gives
\(v_2(5^{2^a}-1)=a+2\).
Indeed squaring an integer \(1+2^b u\), with \(b\geq2\) and \(u\)
odd, increases this valuation by one. Thus \(5\) generates the
\(2^{w-2}\) units which are \(1\bmod4\). Every odd residue is
uniquely \(5^a\) or \(-5^a\). Each of the two cosets contains at
most \(j\) solutions of \(a^j=1\), proving
\eqref{eq:level-roots} in all cases.

Fix now \(Y\geq2\), a prime \(\ell\), and \(j\geq1\).
Let \(S\) be any set of integers \(g\in[Y,2Y]\) satisfying
\(g^j<\ell^{4v_\ell(g^j-1)}\).
If \(S\) is nonempty, put
\[
            w=\min_{g\in S}v_\ell(g^j-1),\qquad Q=\ell^w .
\]
The defining inequality implies \(w\geq1\).
All \(g\in S\) satisfy \(g^j=1\bmod Q\).
For an element \(g_0\) attaining the minimum,
\[
                 Y\leq g_0\leq g_0^j<Q^4 ,
\]
so \(Q>Y^{1/4}\). In an interval of length \(Y\), each residue
class modulo \(Q\) contains at most \(Y/Q+1\) integers.
Equation~\eqref{eq:level-roots} therefore gives
\begin{equation}
 \label{eq:level-exception-fixed}
                 |S|\leq2j(Y^{3/4}+1).
\end{equation}
This is a bound on integers, and hence also on primes.

For any fixed integer \(\ell\), the number of \(g\in[Y,2Y]\) with
\(\ell\in{\cal B}_r(g)\) is at most
\[
 1+\sum_{j=1}^{r}2j(Y^{3/4}+1)
                  \leq1+2r^2(Y^{3/4}+1).
\]
The first term accounts for \(g=\ell\).
Combining this with \eqref{eq:level-bb-card}, the number of
\(g\in[Y,2Y]\) for which
\({\cal B}_r(g)\cap{\cal B}_r(g_1)\ne\varnothing\) is at most
\begin{equation}
 \label{eq:level-excluded}
       (3r+1)\{1+2r^2(Y^{3/4}+1)\}
                       \leq20r^3Y^{3/4}.
\end{equation}
For the last inequality use \(3r+1\leq4r\) and
\(1+2r^2(Y^{3/4}+1)\leq5r^2Y^{3/4}\).

It follows that a set \(G\subset[Y,2Y]\) with
\[
                  |G|\geq Y^{9/10},
       \qquad 20r^3<Y^{3/20}
\]
contains distinct elements \(g_1,g_2\) with disjoint exceptional
sets. Choose any \(g_1\in G\). Formula~\eqref{eq:level-excluded}
excludes strictly fewer than \(|G|\) choices for \(g_2\).
It excludes \(g_1\) itself since \(g_1\in{\cal B}_r(g_1)\), so
the two choices are distinct.

\subsection{The analytic prime input and its scale}

For \(d\equiv3\bmod4\) prime, define
\[
 \begin{split}
 \operatorname{NormOne}(d,N,p)\quad\Longleftrightarrow\quad&
 \exists A,B\in\mathbb Z,\quad A^2+dB^2=4p,\\
 &N\mid B,\qquad 2N\mid A-2-B .
 \end{split}
\]
The short norm-prime estimate used here asserts that there are
absolute constants \(A_0,c,c'>0\) such that, for
\(d\geq100\), \(N\geq3\), \(X\geq1\), and
\[
 \log(4dN^2)\leq c\sqrt{\log X},\qquad
 e^{-c\sqrt{\log X}}\leq h\leq1 ,
\]
the number of primes \(p\in(X,X(1+h)]\) satisfying
\(\operatorname{NormOne}(d,N,p)\) is at least
\begin{equation}
 \label{eq:level-short-norm}
           \frac{c'hX}{(4dN^2)^{A_0}\log X}.
\end{equation}
We will use \eqref{eq:level-short-norm} as a stated input until its
Hecke-family construction and exceptional-zero estimates are
translated. The implication from these family statements is proved below.
Its quantitative weakness is sufficient for all
the following uniform comparisons.
The exact statement is
\srcid{NormPrimesImagQuadShortIntervalLowerBoundStatement};
its subsequent deduction is given in
\srcid{NormPrimesImagQuadShortIntervalLowerBound}.

Suppose
\[
 d\mid N,\qquad
 \log N\leq3x^{\epsilon/5}(\log x)^2,\qquad \epsilon>0 .
\]
For all sufficiently large \(x\), uniformly in \(N,d\) satisfying
these inequalities and all \(Y\) with \(\log Y\geq8x^\epsilon\),
the bound \eqref{eq:level-short-norm} with \(h=1\) gives at least
\(Y^{9/10}\) primes in \((Y,2Y]\).
We spell out uniformity in \(Y\).
Put \(t=x^{\epsilon/5}\), \(L=\log Y\), and \(D=4dN^2\).
For large \(x\),
\((\log x)^2\leq t\), and, since \(d\leq N\),
\[
        \log D\leq\log4+3\log N\leq10t^2,
                  \qquad L\geq8t^5 .
\]
It follows that
\(\log D/\sqrt L\leq(10/\sqrt8)t^{-1/2}\), so the first
hypothesis of \eqref{eq:level-short-norm} holds uniformly.
The remaining size comparison is
\[
 \log\frac{c'Y}{D^{A_0}\log Y}
             -\frac9{10}\log Y
       =\frac L{10}-A_0\log D-\log L+\log c'\geq0 .
\]
The function \(L/10-\log L\) is increasing for \(L\geq10\).
Its minimum over \(L\geq8t^5\) is therefore attained at \(8t^5\),
where
\[
 \frac45t^5-10A_0t^2-\log(8t^5)+\log c'\longrightarrow+\infty .
\]
Thus a threshold depending only on \(\epsilon\) and the absolute
constants works for all allowed \(N,d,Y\).

Take now
\[
             b_x=e^{x^\epsilon},\qquad Y=(k b_x)^8,
             \qquad k,r\geq1,\quad r\leq x^5 .
\]
Then \(\log Y=8(\log k+x^\epsilon)\geq8x^\epsilon\), and
\[
       20r^3\leq20x^{15}<e^{x^\epsilon}
                       \leq(kb_x)^{6/5}=Y^{3/20}
\]
for sufficiently large \(x\). The preceding two subsections supply
two primes \(g_i\in[Y,2Y]\), each satisfying
\(\operatorname{NormOne}(d,N,g_i)\), with
\[
                g_1\ne g_2,\qquad
                  {\cal B}_r(g_1)\cap{\cal B}_r(g_2)=\varnothing .
\]
They also split in the cyclotomic host.
Indeed write \(B=Nb\) and \(A=2+B+2Na\). Substituting in
\(A^2+dB^2=4g_i\) and reducing modulo \(4N\) gives
\(4g_i\equiv4\bmod4N\), hence \(g_i\equiv1\bmod N\).

\subsection{The lower bound from a characterwise explicit formula}

We make the analytic content of \eqref{eq:level-short-norm} more
precise. The following implication accounts for the exceptional real
zero and the removal of prime powers. It can be checked with three
specified inputs about the ray-class family.

Let \(G\) be a finite abelian group of order \(H\), and let \(D\geq4\)
satisfy \(H\leq D\). Suppose nonnegative weights \(\lambda_g(n)\)
satisfy
\[
             \sum_{g\in G}\lambda_g(n)\leq2\Lambda(n).
\]
Write
\(\Psi_\chi(u)=\sum_{n\leq u}\sum_g\chi(g)\lambda_g(n)\).
Assume constants \(a,\eta,C,\gamma>0\), with \(\gamma\leq1\),
have the following uniform properties.
Across all characters there is at most one exceptional pair
\((\chi,\beta)\), where \(3/4\leq\beta<1\) and
\(1-\beta\geq\gamma D^{-a}\).
For \(u\geq2\), \(\log D\leq\sqrt{\log u}\), assume
\begin{equation}
 \label{eq:level-psi-input}
 \left|\Psi_\chi(u)-{\bf1}_{\chi=1}u+
               \sum_{\beta\in Z_\chi}\frac{u^\beta}{\beta}\right|
             \leq Cu e^{-\eta\sqrt{\log u}}.
\end{equation}
Here each \(Z_\chi\) is empty except for the possible exceptional
pair. Finally assume that
\(\lambda_0(p)\ne0\), for \(p\) prime, implies
\(\operatorname{NormOne}(d,N,p)\).

Put \(\delta=\gamma D^{-a}\).
Character orthogonality gives
\[
              \sum_\chi\Psi_\chi(u)
                        =H\sum_{n\leq u}\lambda_0(n).
\]
If \(V=X(1+h)\), \(0<h\leq1\), subtraction of
\eqref{eq:level-psi-input} at \(X\) and \(V\) therefore gives
\begin{equation}
 \label{eq:level-mass}
 H\sum_{X<n\leq V}\lambda_0(n)
 \geq hX-\sum_{\chi,\beta\in Z_\chi}
          \frac{V^\beta-X^\beta}{\beta}
                  -3DCX e^{-\eta\sqrt{\log X}}.
\end{equation}
The factor \(3\) bounds the two endpoint errors since \(V\leq2X\).

If an exceptional zero exists and \(\log X\geq1\), then
\[
 \frac{V^\beta-X^\beta}{\beta}
       \leq hX^\beta
       =hX e^{-(1-\beta)\log X}
       \leq hX e^{-\delta}
       \leq hX(1-\delta/2).
\]
The first inequality follows from the concavity of
\(t^\beta\) at \(1\). For the last one, use
\(e^{-\delta}\leq1/(1+\delta)\leq1-\delta/2\), valid for
\(0\leq\delta\leq1\). If there is no exceptional zero, the sum
in \eqref{eq:level-mass} is zero, and the same bound on that sum
still applies.

Choose \(c>0\) small enough that \(c\leq1\) and
\((a+2)c<\eta\). By decreasing \(c\) further, all
\(\log D\leq cL\), \(L=\sqrt{\log X}\), imply
\[
       3DCe^{-\eta L}
                     \leq\frac14 e^{-cL}\gamma D^{-a}.
\]
To justify the uniform choice, the ratio of the left side to the
right side is at most
\[
       \frac{12C}{\gamma}
             e^{-\{\eta-(a+2)c\}L}.
\]
Also \(L\geq\log4/c\). Take \(c\) sufficiently small that the last
display is at most one at this lower endpoint. It is decreasing
in \(L\). Hence, when \(h\geq e^{-cL}\),
\eqref{eq:level-mass} implies
\[
          \sum_{X<n\leq V}\lambda_0(n)
                       \geq\frac{\gamma hX}{4D^{a+1}}.
\]

Let \(P\) count the norm primes in \((X,V]\).
The prime terms in the last sum are at most \(2\log V\cdot P\).
The other terms are at most
\(2\{\psi(V)-\theta(V)\}\).
An elementary bound sufficient here is
\[
 \psi(V)-\theta(V)
       =\sum_{j=2}^{\lfloor\log V/\log2\rfloor}
                         \theta(V^{1/j})
       \leq\frac{\sqrt V(\log V)^2}{\log2}.
\]
Indeed there are at most \(\log V/\log2\) terms, and each is at
most \(\sqrt V\log V\), by counting all integers up to
\(V^{1/j}\). Since \(V\leq2X\), \(\log V\leq2\log X\) for
\(X\geq2\), the unwanted contribution is at most
\(C_1\sqrt X(\log X)^2\) for an absolute \(C_1\).
Its ratio to \(\gamma hX/(8D^{a+1})\) is bounded by
\[
 \frac{8C_1}{\gamma}L^4
        \exp\{-L^2/2+(a+2)cL\}.
\]
Decreasing \(c\) again makes this at most one for all
\(L\geq\log4/c\), because the negative quadratic term dominates
the linear term and \(\log L\). Thus at least
\(\gamma hX/(8D^{a+1})\) of the mass comes from primes.
Dividing by \(2\log V\leq4\log X\) gives
\[
             P\geq
               \frac{\gamma hX}{32D^{a+1}\log X}.
\]
This is \eqref{eq:level-short-norm} with
\(A_0=a+1\), \(c'=\gamma/32\), and \(D=4dN^2\).

The analytic inputs are now explicit.
The ray-class construction must supply the nonnegative weights and
their norm-prime support. The family explicit formula must give
\eqref{eq:level-psi-input} with a unique possible exceptional zero.
Section~\ref{sec:ray-real-gap} proves the real value and zero-gap estimates
from the counting, theta and positive-convolution identities of the family. In the release these are
\srcid{ExistsLFamRayStatement},
\srcid{RayClassFamilyPsiWithExceptionalZeroStatement}, and
\srcid{ExistsRayClassLFunctionRealZeroFarFromOneStatement}.
The last is itself obtained by combining a polynomial lower bound
for \(|F_\chi(1)|\) with
\[
             |F'_\chi(z)|\leq18000D
                           \qquad(|z-1|\leq1/4).
\]
If \(F_\chi(\beta)=0\), integration on \([\beta,1]\) gives
\(|F_\chi(1)|\leq18000D(1-\beta)\).
A bound \(c_0D^{-a_0}\leq|F_\chi(1)|\) therefore supplies
\(1-\beta\geq(c_0/18000)D^{-(a_0+1)}\).
This identifies the remaining analytic statements without suppressing
their dependence on \(D\).

\subsection{Three distinct small splitting primes}

For the second input, Linnik's theorem supplies a constant \(A\geq1\)
such that for every
integer \(M\geq2\) there is a prime \(p\equiv1\bmod M\) with
\(p\leq M^A\). Its formal name is \srcid{SmallOne}. For example, Xylouris's
Theorem~1.1 gives \(p\leq C M^{5.18}\) uniformly for each reduced
residue class. Replacing the exponent by
\(A=5.18+\max(0,\log C/\log2)\) absorbs the absolute constant for
all \(M\geq2\). This supplies the precise version used here.
\footnote{T.~Xylouris, \emph{On the least prime in an arithmetic
progression and estimates for the zeros of Dirichlet \(L\)-functions},
Acta Arithmetica 150 (2011), Theorem~1.1,
\href{https://doi.org/10.4064/aa150-1-4}{doi:10.4064/aa150-1-4}.}
For \(N\geq30\), first apply it to \(15N\), then to \(15Np\),
and finally to \(Npp'\). The resulting primes satisfy
\[
 \begin{gathered}
 p,p'\equiv1\bmod15N,\qquad q_0\equiv1\bmod N,\\
 p'\equiv1\bmod p,\qquad q_0\equiv1\bmod pp' .
 \end{gathered}
\]
They are pairwise distinct, as a prime cannot be both zero and one
modulo itself. Since \(15N\leq N^2\),
\[
 \begin{split}
 p&\leq N^{2A},\\
 p'&\leq (15Np)^A\leq N^{(2A+2)A},\\
 q_0&\leq (Npp')^A
             \leq N^{A\{1+2A+(2A+2)A\}}.
 \end{split}
\]
The common exponent
\[
                    C_8=A\{1+2A+(2A+2)A\}
\]
dominates all three. With \(N\) on the window above,
\[
 \log N^{C_8}
     \leq3C_8x^{\epsilon/5}(\log x)^2<x^\epsilon
\]
for all sufficiently large \(x\). Thus \(p,p',q_0<b_x\), whereas
\(g_i\geq(kb_x)^8\geq b_x\).
In particular \(p,p'\) differ from both \(g_i\), as required.

Every threshold used here depends on the fixed \(\epsilon\) and the
absolute analytic constants. They can be combined by taking their
maximum. Consequently the same choices of \(x,N,d,k,r\) are retained
throughout the construction. The remaining analytic dependencies are the construction of the ray-class
family with the stated identities and its characterwise explicit formula. The small splitting primes follow from the
stated classical least-prime theorem.

\paragraph{Source map.}
The definitions of the exceptional sets are in
\srcid{Formalism.CyclotomicNorms.StatementI}. The counting proofs
are \srcid{ExceptionalPrimesFiniteCardLe},
\srcid{CardPowEqOneLeOfIsCyclicUnits}, and
\srcid{DivLeRpowThreeQuarters}. Their use on the scale is
\srcid{ShortNorms.normPrimesN} and \srcid{ShortNorms.typePrimesN}.
The three-prime construction and its simultaneous application are
\srcid{ShortNorms.ExistsThreePrimesOneModLeRpow} and
\srcid{ShortNorms.levelData}.
\footnote{\href{https://github.com/Deicyde/riemann-hypothesis/tree/40e56bab026552e089d7afcfeff08d2defb9c005/src/Development/AutomorphicArithmetic/ShortNorms}{The prime-selection modules};
\href{https://github.com/Deicyde/riemann-hypothesis/blob/40e56bab026552e089d7afcfeff08d2defb9c005/src/Development/AbelianVarietyTheory/Proofs/FoundationProofs/DivLeRpowThreeQuarters.lean}{the fixed-prime residue count}.}
